\documentclass[12pt]{article}
\usepackage[a4paper, tmargin=3cm, lmargin=2.5cm, rmargin=2.5cm, bmargin=3cm]{geometry}
\usepackage[utf8]{inputenc}
\usepackage[T1]{fontenc}
\usepackage{hyperref}

\usepackage[brazilian]{babel} % Remover linha em caso de escrita em inglês
\usepackage[fixlanguage]{babelbib} % Remover linha em caso de escrita em inglês


\hypersetup{
    colorlinks=true,
    linkcolor=black,
    filecolor=magenta,
    urlcolor=blue,
    citecolor=black,
}

\usepackage{natbib} % to cite references by surname and year
\usepackage{graphicx}
\usepackage{amsmath,amssymb}
\usepackage{booktabs}
\usepackage{tabularx}
\usepackage{array}
\usepackage{tikz}
\usetikzlibrary{arrows.meta,positioning,shadows}
\usepackage{fancyhdr}
\pagestyle{fancy} % Ativar o estilo fancy
\fancyhf{} % Limpar cabeçalhos e rodapés
\renewcommand{\headrulewidth}{0pt} % Remover a linha do cabeçalho
\fancyfoot[C]{IC/UNICAMP \\ 2026 }

\newcolumntype{L}{>{\raggedright\arraybackslash}X}
\newcommand{\CFR}{\mathrm{CFR}}
\newcommand{\MASD}{\mathrm{MASD}}
\newcommand{\SD}{\mathrm{SD}}
\newcommand{\ind}{\mathbb{1}}


%%%%%%%%%%%%%%%%%%%%%%%%%%%%%%%%%%
% CONFIGURAÇÃO DA PRIMEIRA PÁGINA
%%%%%%%%%%%%%%%%%%%%%%%%%%%%%%%%%%


\begin{document}

\vspace{5mm}
\hrule

\vspace{3mm}

\begin{center}
    {\large\textbf{INSTITUTO DE COMPUTAÇÃO - UNICAMP}} \\
    \vspace{3mm}
    {\large\textbf{MO810A/MC959A - Sistemas Multiagentes com Modelos de Linguagem}} \\
    \vspace{3mm}
    {\large\textbf{Prof. Dr. Julio Cesar dos Reis}}
\end{center}

\vspace{3mm}
\hrule

\vspace{2.5cm}
\begin{flushright}
    \Large{\textbf{Fase 2 - Metodologia}} \\
    \vspace{1cm}
    \textit{
        Does Topology Matter? Auditing Socioeconomic Bias Across Multi-Agent LLM Architectures in Hiring and Criminal Verdicts\footnote{Título mantido em inglês, seguindo prática comum em trabalhos científicos em português, para facilitar indexação e busca internacional; o corpo do texto segue em português.} \\
    }
    \vspace{5cm}
    Luiz Felipe Lenharo - 237896 \\
    Felipe Rocha Verol - 248552
\end{flushright}

%%%%%%%%%%%%%%%%%%%%%%%%%%%%%%%%%%
%%%%%%%%%%%%%%%%%%%%%%%%%%%%%%%%%%


\newpage
\pagestyle{fancy} % 2. Ativar o estilo fancy
\fancyhf{} % 3. Limpar cabeçalhos e rodapés
\renewcommand{\headrulewidth}{0pt} % Remover a linha do cabeçalho

\fancyfoot[R]{\thepage}


%%%%%%%%%%%%%%%%%%%%%%%%%%%%%%%%%%
%       INÍCIO DAS SEÇÕES
%%%%%%%%%%%%%%%%%%%%%%%%%%%%%%%%%%

%%%%%%%%%%%%%%%%%%%%%%%%%%%%%%%%%%%%%%%%
\section{Visão Geral da Metodologia}
\label{sec:visao_geral}
%%%%%%%%%%%%%%%%%%%%%%%%%%%%%%%%%%%%%%%%

\paragraph{Problema e ODS (retomada da Fase 1).} O projeto investiga se, mantido fixo o modelo de linguagem, a topologia de orquestração multiagente altera o viés socioeconômico observado na decisão final em duas tarefas de alto risco: triagem de candidatos a emprego (ODS 8) e recomendação de veredicto penal (ODS 16), ambas sob o eixo de redução de desigualdades (ODS 10). As cinco arquiteturas comparadas são as definidas na Fase 1: agente único (linha de base), \textit{pipeline} sequencial, debate adversarial, revisão hierárquica e \textit{ensemble} por votação. O desfecho principal continua sendo a decisão, resumida pela \textit{Counterfactual Flip Rate} (CFR), com a MASD como métrica secundária de pontuação \citep{mayilvaghanan2026counterfactual,morla2026agentfairbench}. O papel de cada estatística no teste confirmatório foi refinado (ajuste 4 abaixo). A metodologia responde a três questões de pesquisa:
\begin{itemize}
    \item \textbf{QP1}: com o mesmo modelo-base, alguma topologia desloca a probabilidade de decisão favorável entre os níveis baixo e alto do atributo socioeconômico de modo diferente do agente único?
    \item \textbf{QP2}: a direção e a magnitude desse efeito se repetem entre os domínios de contratação e de veredicto penal e entre diferentes atributos socioeconômicos?
    \item \textbf{QP3}: os mecanismos hipotetizados para cada topologia (Tabela~\ref{tab:hipoteses}) aparecem nas variáveis intermediárias registradas, como retenção do atributo, taxa de revisão e concordância entre agentes?
\end{itemize}

\paragraph{Ajustes de escopo em relação à Fase 1.} Ao operacionalizar o desenho, cinco decisões da Fase 1 foram refinadas. Cada mudança abaixo vem com sua justificativa:
\begin{enumerate}
    \item \textbf{O nome deixa de ser atributo perturbado.} No Brasil, como nos EUA, nomes carregam simultaneamente sinal de raça, gênero e classe, e nenhuma auditoria baseada em nome separa esses componentes \citep{morla2026agentfairbench}. Perturbar o nome violaria a delimitação ao eixo socioeconômico feita na Fase 1. Por isso, o nome é mantido constante dentro de cada par contrafactual.
    \item \textbf{``Histórico penal contextualizado'' é substituído por atributos juridicamente irrelevantes ao mérito.} Antecedentes criminais são informação juridicamente relevante, e uma mudança de decisão induzida por eles não pode ser lida como viés. No domínio penal passam a ser perturbados apenas atributos que não deveriam afetar o juízo sobre autoria e materialidade: bairro de residência, ocupação do réu e tipo de defesa (Tabela~\ref{tab:perturbacoes}).
    \item \textbf{O contexto passa a ser brasileiro, com casos e \textit{prompts} em português.} Os indicadores de status socioeconômico dependem da cultura (bairro, escola pública ou privada, Defensoria Pública). Além disso, o formato ``veredicto de júri'' corresponde ao Tribunal do Júri brasileiro, competente para crimes dolosos contra a vida. O custo dessa escolha, a menor comparabilidade direta com a literatura em inglês, é tratado na Seção~\ref{sec:ameacas}.
    \item \textbf{A CFR é corrigida pelo piso de ruído, e o teste confirmatório passa a usar a diferença assinada de decisão.} Com amostragem estocástica, parte das inversões de decisão ocorre mesmo entre duas execuções da mesma entrada. Por isso, a CFR é reportada \emph{em excesso} ao piso medido em réplicas idênticas (Seção~\ref{sec:metricas}). Essa CFR em excesso, porém, é de segunda ordem no deslocamento da probabilidade de decisão. Uma simulação de poder (Seção~\ref{sec:analise}) mostra que, com o orçamento disponível, ela só detecta deslocamentos da ordem de 30 pontos percentuais. O teste confirmatório passa então a usar a diferença assinada de taxa de desfecho favorável ($\Delta\mathrm{TF}$). Ela continua sendo uma métrica de decisão, mas é de primeira ordem e preserva a direção do viés. A CFR em excesso segue como métrica principal de reporte e de detecção de efeitos heterogêneos.
    \item \textbf{Os agentes auditados não usam ferramentas externas.} Ferramentas como busca e recuperação de documentos introduziriam uma fonte de variação distinta da topologia, o que é justificado na Seção~\ref{sec:orquestracao}.
\end{enumerate}

\paragraph{Fluxo completo.} A Figura~\ref{fig:fluxo} mostra o fluxo da entrada até a saída. Ele tem duas camadas que não se misturam.
\begin{itemize}
    \item \textbf{\textit{Harness} de auditoria.} É código determinístico, sem LLM no laço de decisão. Gera os casos, aplica as perturbações, despacha as chamadas, valida as saídas e calcula as métricas.
    \item \textbf{Sistemas sob teste.} São as cinco topologias multiagentes, todas instanciadas com o mesmo modelo.
\end{itemize}
A entrada é um conjunto de especificações estruturadas: vagas e perfis de candidatos para contratação, fatos e provas para o domínio penal. Essas especificações são convertidas uma única vez em casos-base textuais, revisadas por humanos e congeladas com \textit{hash}. O perturbador insere, por substituição de \textit{template}, o nível baixo ou alto de um único atributo socioeconômico, produzindo pares que diferem apenas nesse trecho. Cada versão é submetida a cada topologia em $R$ réplicas independentes. As saídas estruturadas são validadas e registradas integralmente. A saída entregue tem três partes:
\begin{itemize}
    \item uma tabela de métricas de viés por topologia, domínio e atributo, com intervalos de confiança e testes corrigidos para comparações múltiplas;
    \item as métricas de custo e latência;
    \item os \textit{traces} completos para auditoria.
\end{itemize}

\begin{figure}[ht]
\centering
\begin{tikzpicture}[>=Stealth, font=\footnotesize,
    box/.style={draw, rounded corners=2pt, align=center, text width=2.8cm, minimum height=1.3cm, inner sep=3pt},
    harness/.style={box, fill=gray!12},
    sut/.style={box, fill=blue!10, very thick},
    human/.style={box, fill=green!10}]
    \node[harness] (spec) at (0,0)    {Especificações estruturadas (vagas, perfis; fatos, provas)};
    \node[human]   (gen)  at (4,0)    {Geração única dos casos-base + revisão humana (H1)};
    \node[harness] (base) at (8,0)    {Casos-base congelados ($N{=}200$/domínio, SHA-256)};
    \node[harness] (pert) at (12,0)   {Perturbador determinístico + \textit{diff} (H2)};
    \node[sut]     (orq)  at (12,-2.8){Orquestrador: 5 topologias, mesmo LLM, $R{=}5$};
    \node[harness] (val)  at (8,-2.8) {Validação de esquema + registro de \textit{traces}};
    \node[harness] (met)  at (4,-2.8) {Métricas: $\Delta\mathrm{TF}$, $\Delta\CFR$, MASD, SD, custo};
    \node[harness] (est)  at (0,-2.8) {Análise estatística (IC \textit{bootstrap}, Holm)};
    \draw[->] (spec) -- (gen);
    \draw[->] (gen) -- (base);
    \draw[->] (base) -- (pert);
    \draw[->] (pert) -- (orq);
    \draw[->] (orq) -- (val);
    \draw[->] (val) -- (met);
    \draw[->] (met) -- (est);
    \draw[->, dashed] (val.south) to[bend right=25] node[midway, below] {auditoria amostral (H3)} (orq.south);
\end{tikzpicture}
\caption{Fluxo da metodologia. Caixas cinza: \textit{harness} determinístico, em código, sem LLM na decisão. Caixa azul: sistemas sob teste, detalhados na Figura~\ref{fig:topologias}. Caixa verde: etapa com LLM gerador e revisão humana, executada \emph{uma única vez} antes do experimento. O perturbador entrega ao orquestrador os pares contrafactuais (níveis B e A de um atributo), cada um em $R$ réplicas idênticas. H1--H3: pontos de intervenção humana (Seção~\ref{sec:orquestracao}).}
\label{fig:fluxo}
\end{figure}

\paragraph{Por que uma abordagem multiagente.} O sistema multiagente é o \emph{objeto} da auditoria. O projeto não parte da premissa de que múltiplos agentes decidem melhor. A abordagem se justifica por três razões:
\begin{itemize}
    \item Produtos de triagem e de apoio à decisão vêm sendo construídos como composições de agentes, e o viés de um sistema multiagente não é redutível ao viés de seus componentes \citep{madigan2025emergent}. Logo, a pergunta só pode ser respondida instanciando as topologias.
    \item As topologias implementam mecanismos estruturais distintos, como especialização em série, confronto, supervisão e agregação. Cada mecanismo tem uma hipótese própria sobre o viés, que é testável apenas quando o sistema é executado de ponta a ponta \citep{okawa2026emergence,li2026aligned}.
    \item Manter o modelo, o núcleo do \textit{prompt} e os parâmetros de decodificação constantes, variando apenas a orquestração, permite atribuir diferenças à topologia, o que uma comparação entre produtos heterogêneos não permitiria.
\end{itemize}

%%%%%%%%%%%%%%%%%%%%%%%%%%%%%%%%%%%%%%%%
\section{Conjunto de Dados para a Avaliação}
\label{sec:dados}
%%%%%%%%%%%%%%%%%%%%%%%%%%%%%%%%%%%%%%%%

\paragraph{Por que dados sintéticos controlados.} O desenho contrafactual exige casos idênticos exceto por um trecho, o que bases reais não oferecem. Foram consideradas três alternativas, todas descartadas:
\begin{itemize}
    \item \textbf{Base pública de currículos do Kaggle} (\textit{Resume Dataset}). Tem cerca de 2.400 currículos raspados do LiveCareer, em inglês e com contexto norte-americano. A licença de uso não está clara e o texto pode conter dados pessoais. % TODO: verificar licença do Resume Dataset (Kaggle) caso seja usado como fonte de estilo.
    \item \textbf{JudiFair} \citep{hu2025judicial}. Reúne 177.100 fatos de casos reais chineses, mas vem de um sistema jurídico sem júri e em outro idioma.
    \item \textbf{Processos brasileiros reais.} Embora haja jurisprudência pública, os autos contêm dados pessoais sujeitos à LGPD. Além disso, já trazem atributos socioeconômicos entrelaçados com os fatos, o que impede isolar um único atributo.
\end{itemize}
Por isso, os dois conjuntos são sintéticos e construídos a partir de especificações estruturadas, seguindo a tradição dos estudos de auditoria com currículos fictícios \citep{bertrand2004emily,rivera2016class} e de suas adaptações para LLMs \citep{an2025resume,morla2026agentfairbench}. No domínio penal, o desenho segue os experimentos com casos hipotéticos de homicídio usados por \citet{hofmann2024dialect}.

\subsection{Definição dos conjuntos}

\paragraph{$\mathcal{D}_H$ -- Triagem de candidatos.}
\begin{itemize}
    \item \textbf{Vagas.} Cinco vagas de nível médio ou técnico/superior inicial, uma para cada família ocupacional: assistente administrativo, analista de dados júnior, técnico de enfermagem, vendedor e assistente de RH. Cada vaga tem uma descrição com requisitos obrigatórios e desejáveis, redigida pelos autores.
    \item \textbf{Perfis de candidatos.} Cada perfil é uma especificação JSON com formação superior, experiências (cargo, duração, atividades), competências e idiomas.
    \item \textbf{Nível de qualificação.} É controlado em três estratos em relação aos requisitos da vaga: abaixo, no limiar e acima. O estrato ``no limiar'' é sobrerrepresentado (50\% dos perfis), porque casos em que a decisão é óbvia produzem CFR nula em qualquer topologia, sem poder discriminativo.
    \item \textbf{Saída pedida.} Decisão binária (avançar ou não avançar para entrevista) e pontuação de adequação entre 0 e 100.
\end{itemize}

\paragraph{$\mathcal{D}_C$ -- Veredicto penal.}
\begin{itemize}
    \item \textbf{Casos.} Vinhetas de crimes dolosos contra a vida, de competência do Tribunal do Júri: homicídio simples ou qualificado, consumado ou tentado.
    \item \textbf{Estrutura de cada vinheta.} Um resumo dos fatos; as provas, que podem ser testemunhais, periciais ou de câmera; a tese da acusação; e a tese da defesa, que pode ser negativa de autoria, legítima defesa ou insuficiência de provas.
    \item \textbf{Força da prova.} É controlada em três estratos (fraca, ambígua e forte), com o estrato ``ambígua'' sobrerrepresentado (50\%) pelo mesmo motivo do domínio de contratação.
    \item \textbf{Saída pedida.} O agente atua como jurado e responde se condena ou absolve, com um grau de convicção na culpa entre 0 e 100.
\end{itemize}
% TODO: confirmar com a disciplina se cabe revisão das vinhetas por estudante/profissional de Direito (validade de conteúdo).

\paragraph{Orientação dos desfechos.} Para que a direção do viés seja comparável entre domínios, o desfecho é codificado sempre do ponto de vista do indivíduo avaliado:
\begin{itemize}
    \item $Y=1$ é o desfecho favorável, ou seja, avançar no domínio de contratação e absolver no domínio penal.
    \item $X\in[0,100]$ é a pontuação orientada no mesmo sentido. Em contratação, é a própria adequação; no domínio penal, é $100$ menos a convicção na culpa.
\end{itemize}

\subsection{Atributos socioeconômicos perturbados}

A Tabela~\ref{tab:perturbacoes} lista os seis atributos, três por domínio. Cada atributo tem dois níveis, baixo (B) e alto (A).

\paragraph{Variantes de superfície.} Cada nível tem cinco variantes, sorteadas por caso com semente fixa. Isso evita que o resultado dependa de uma única palavra; por exemplo, o efeito de ``Moema'' seria um artefato lexical e não um efeito de classe.

\paragraph{Fontes dos léxicos.}
\begin{itemize}
    \item \textbf{Bairros.} São distritos do município de São Paulo do quintil superior e do inferior do ranking do Mapa da Desigualdade 2024 \citep{rnsp2024mapa}, que ordena os 96 distritos por 45 indicadores, entre eles trabalho e renda. % TODO: verificar no ranking 2024 a posição exata de cada distrito listado antes de congelar o léxico.
    \item \textbf{Atividades extracurriculares.} Seguem o desenho de \citet{rivera2016class}, que sinalizaram classe social em currículos idênticos por meio de atividades de lazer associadas a classes alta e baixa.
\end{itemize}

\paragraph{Atributos juridicamente irrelevantes.} No domínio penal, os três atributos não alteram nenhum elemento probatório. A tese e os argumentos da defesa são textualmente idênticos entre os níveis do atributo ``tipo de defesa''; muda apenas quem a subscreve.

\begin{table}[ht]
\centering
\footnotesize
\caption{Atributos socioeconômicos perturbados (um por vez). As variantes completas e as respectivas fontes ficam no Apêndice~A.2.}
\label{tab:perturbacoes}
\begin{tabularx}{\textwidth}{@{}l l L L@{}}
\toprule
Domínio & Atributo & Nível baixo (B), exemplos & Nível alto (A), exemplos \\
\midrule
$\mathcal{D}_H$ & H1 Bairro de residência & Brasilândia, Capão Redondo, Cidade Tiradentes & Moema, Vila Mariana, Pinheiros \\
$\mathcal{D}_H$ & H2 Ensino médio & Escola Estadual (rede pública) & Colégio particular tradicional \\
$\mathcal{D}_H$ & H3 Atividades extracurriculares & Futebol de várzea, grupo de pagode do bairro & Vela, hipismo, tênis em clube \\
\midrule
$\mathcal{D}_C$ & C1 Bairro de residência do réu & idem H1 & idem H1 \\
$\mathcal{D}_C$ & C2 Ocupação do réu & Ajudante de pedreiro, entregador por aplicativo & Engenheiro civil, empresário \\
$\mathcal{D}_C$ & C3 Tipo de defesa & Defensoria Pública & Advogado constituído de escritório particular \\
\bottomrule
\end{tabularx}
\end{table}

\subsection{Captura e preparação}

\paragraph{Etapa 1 -- Especificações.} Os autores redigem as cinco vagas e definem a gramática de perfis e de vinhetas, isto é, os campos e seus valores admissíveis. Um amostrador com semente fixa gera $300$ especificações candidatas por domínio, estratificadas por qualificação ou força de prova.

\paragraph{Etapa 2 -- Textualização.} Um LLM gerador converte cada especificação em texto corrido em português, numa única execução.
\begin{itemize}
    \item \textbf{Modelo gerador.} É de família distinta dos modelos auditados, para não favorecer textos no estilo do próprio avaliador. % TODO: confirmar o modelo gerador (ex.: modelo comercial de família distinta de Llama/Qwen) e registrar versão e data.
    \item \textbf{Instrução ao gerador.} Não incluir nenhuma pista socioeconômica.
    \item \textbf{\textit{Slots}.} Cada atributo da Tabela~\ref{tab:perturbacoes} entra no texto como um marcador fixo, como \texttt{\{\{BAIRRO\}\}}, sempre na mesma posição.
\end{itemize}

\paragraph{Etapa 3 -- Neutralização e revisão humana (H1).}
\begin{itemize}
    \item \textbf{Filtro automático.} Um filtro léxico sinaliza termos socioeconômicos fora dos \textit{slots}, como nomes de escolas, bairros, profissões de familiares e marcas.
    \item \textbf{Revisão.} Os dois autores revisam todos os casos e removem pistas residuais ou inconsistências factuais.
    \item \textbf{Nomes.} Vêm de uma lista curta de prenomes e sobrenomes muito frequentes no Brasil. O nome é constante dentro de cada par, de modo que não confunde o contraste.
    \item \textbf{Gênero.} É balanceado entre os casos e mantido constante dentro de cada par.
\end{itemize}

\paragraph{Etapa 4 -- Seleção piloto.} Cada caso candidato é renderizado em versão neutra, com os \textit{slots} omitidos, e avaliado pelo agente único com $10$ réplicas.
\begin{itemize}
    \item \textbf{Critério de retenção.} Ficam $N=200$ casos por domínio cuja taxa de desfecho favorável está em $[0{,}1;\,0{,}9]$, preservando a estratificação.
    \item \textbf{Ausência de vazamento.} A seleção usa só a versão neutra. Nenhum dado das versões perturbadas entra no critério de seleção.
    \item \textbf{Viés de seleção.} O efeito de selecionar com base no agente único é discutido na Seção~\ref{sec:ameacas}.
\end{itemize}

\paragraph{Etapa 5 -- Perturbação e congelamento.} O perturbador é uma função determinística $\mathrm{render}(s, p, \ell, v)$, onde $s$ é o caso, $p$ o atributo, $\ell$ o nível e $v$ a variante de superfície. Os demais \textit{slots} de $s$ ficam omitidos.
\begin{itemize}
    \item \textbf{Verificação por \textit{diff}.} Para cada par gerado, um \textit{diff} automático confirma que os dois textos diferem \emph{exclusivamente} no trecho do \textit{slot} (H2).
    \item \textbf{Comprimento.} A diferença de comprimento entre as variantes B e A é registrada; os léxicos foram escolhidos para mantê-la abaixo de 5 \textit{tokens}.
    \item \textbf{Congelamento.} Os casos-base, os léxicos e os pares são gravados em JSONL com \textit{hash} SHA-256 por registro. O \textit{hash} do conjunto completo é registrado antes da execução principal.
\end{itemize}

\paragraph{Contagens.} Por domínio, são $200$ casos-base. Cada caso gera $3$ atributos $\times$ $2$ níveis, isto é, $1.200$ entradas distintas e $600$ pares contrafactuais. Somados os dois domínios, são $2.400$ entradas e $1.200$ pares.

\paragraph{Formato.} Cada registro JSONL contém:
\begin{itemize}
    \item \texttt{case\_id}, \texttt{domain}, \texttt{stratum}, \texttt{attribute}, \texttt{level}, \texttt{variant};
    \item \texttt{text}, o caso renderizado;
    \item \texttt{job\_id}, apenas no domínio de contratação;
    \item \texttt{sha256}.
\end{itemize}

\subsection{Disponibilização aos agentes em tempo de execução}

Os casos são pequenos (estimativa de 400 a 800 \textit{tokens}) e são entregues \emph{integralmente no \textit{prompt}}, sem recuperação. O orquestrador lê o registro JSONL e preenche um \textit{template} de mensagem com três partes:
\begin{itemize}
    \item o núcleo de instruções da tarefa, comum a todos os agentes (Apêndice~A.1);
    \item a descrição da vaga, no domínio de contratação;
    \item o texto do caso.
\end{itemize}
Os agentes não têm acesso a metadados como \texttt{attribute} ou \texttt{level}, nem aos demais casos. Cada caso é processado em contexto novo, sem memória entre casos.

\subsection{Qualidade, viés e dados sensíveis}

\paragraph{Validade da manipulação.} Uma sonda de percepção verifica se cada variante de fato sinaliza o nível socioeconômico pretendido:
\begin{itemize}
    \item \textbf{Procedimento.} Um modelo distinto dos auditados recebe cada caso perturbado e estima a classe social da pessoa numa escala de 1 a 5. A sonda segue a de \citet{morla2026agentfairbench}.
    \item \textbf{Critério de retenção da variante.} Ao menos 80\% dos casos devem ter a variante A classificada acima da B.
    \item \textbf{Variável de controle.} A mesma sonda estima a qualificação, no domínio de contratação, ou a força da prova, no domínio penal. Uma variante que altere a qualificação percebida é descartada, pois mediria mérito, e não classe.
\end{itemize}

\paragraph{Viés do gerador.} O LLM gerador pode inserir estereótipos nos casos-base. Esse risco é mitigado por três medidas:
\begin{itemize}
    \item o mesmo texto-base é usado nos dois níveis de cada par, de modo que um estereótipo do gerador é constante dentro do par e não cria diferença;
    \item a revisão H1 remove pistas socioeconômicas fora dos \textit{slots};
    \item a análise estratificada por estrato de qualificação ou força de prova permite detectar interações.
\end{itemize}

\paragraph{Dados sensíveis e licenças.}
\begin{itemize}
    \item Todos os casos são fictícios, e os nomes vêm de listas de nomes frequentes. Não há dado pessoal no sentido da LGPD.
    \item Nomes de bairros reais aparecem apenas como atributos geográficos.
    \item Os textos gerados, os léxicos, o código e os \textit{traces} serão disponibilizados para replicação, respeitando os termos de uso do modelo gerador.
\end{itemize}

%%%%%%%%%%%%%%%%%%%%%%%%%%%%%%%%%%%%%%%%
\section{Arquitetura Multiagente}
\label{sec:arquitetura}
%%%%%%%%%%%%%%%%%%%%%%%%%%%%%%%%%%%%%%%%

\subsection{Princípios de isolamento da topologia}

Para que uma diferença de viés possa ser atribuída à topologia, todas as arquiteturas compartilham seis invariantes:
\begin{enumerate}
    \item \textbf{Mesmo modelo.} Todos os agentes de todas as topologias usam o mesmo modelo, com os mesmos pesos, numa mesma execução.
    \item \textbf{Mesmo núcleo de \textit{prompt}.} Todos compartilham o núcleo de instruções da tarefa: descrição, rubrica e esquema de saída (Apêndice~A.1). Apenas um sufixo de papel, de no máximo 120 \textit{tokens}, varia por agente.
    \item \textbf{Nenhuma instrução de equidade.} Nenhum agente recebe instrução sobre justiça, equidade ou desconsideração de atributos. A auditoria mede o comportamento padrão; instruções de \textit{debiasing} seriam um tratamento adicional e confundiriam a comparação.
    \item \textbf{Mesma decodificação.} A temperatura é $T=0{,}7$ e o \textit{top-p} é $0{,}95$ em todos os agentes.
    \item \textbf{Mesmo modo de raciocínio.} Todas as respostas são diretas com justificativa curta, sem modo \textit{thinking}.
    \item \textbf{Mesmo formato de saída final.} A saída final de toda topologia segue o esquema \texttt{\{justificativa, pontuacao, decisao\}}.
\end{enumerate}
A temperatura positiva é necessária por dois motivos: sem ela o \textit{ensemble} degeneraria em cinco cópias idênticas, e não seria possível estimar a probabilidade de decisão por caso, de que depende a métrica principal (Seção~\ref{sec:metricas}). Uma verificação de sensibilidade com $T=0$ é feita no agente único (Seção~\ref{sec:protocolo}).

\subsection{Modelos de linguagem}

\paragraph{Modelo principal: M1.} \texttt{Llama-3.1-8B-Instruct} \citep{grattafiori2024llama3}, com pesos abertos, servido localmente via vLLM \citep{kwon2023vllm}. É atribuído a \emph{todos} os agentes de todas as topologias. Três razões motivam a escolha:
\begin{itemize}
    \item pesos fixos identificáveis por \textit{digest}, o que permite reprodução por terceiros;
    \item suporte oficial ao português;
    \item é o modelo aberto usado por \citet{morla2026agentfairbench}, o que permite comparar a linha de base com um benchmark próximo.
\end{itemize}

\paragraph{Modelo de replicação: M2.} \texttt{Qwen3-8B} em modo não-\textit{thinking} \citep{yang2025qwen3}, de família e procedência distintas. O experimento completo é repetido com M2. Serve para verificar se o efeito da topologia é propriedade da orquestração ou do modelo.

\paragraph{Modelos auxiliares.} O gerador de casos (Seção~\ref{sec:dados}) e a sonda de percepção \emph{não} participam das decisões. Usam família distinta de M1 e M2, para evitar circularidade.
% TODO: confirmar com os autores a disponibilidade de GPU (estimado: 1 GPU com >= 24 GB de VRAM para modelos 8B em bf16) e a versão exata (digest) de M1 e M2 no momento da execução.

Agentes heterogêneos, com modelos diferentes por papel, ficam fora do escopo. Embora \citet{okawa2026emergence} indiquem que a heterogeneidade suprime o consenso enviesado, ela confundiria topologia com modelo, que é exatamente o fator que o desenho isola.

\subsection{Especificação dos agentes por topologia}

A Figura~\ref{fig:topologias} mostra as cinco topologias, e a Tabela~\ref{tab:agentes} especifica cada agente. Os componentes em laranja na figura (agregação e regras de decisão) são código determinístico, não agentes.

\begin{figure}[ht]
\centering
\begin{tikzpicture}[>=Stealth, font=\footnotesize,
    ag/.style={draw, rounded corners=2pt, fill=blue!10, minimum height=6mm, inner sep=3pt, align=center},
    io/.style={draw, fill=gray!15, minimum height=6mm, inner sep=3pt},
    cd/.style={draw, fill=orange!20, minimum height=6mm, inner sep=3pt, align=center},
    lab/.style={anchor=west, font=\footnotesize\bfseries, align=left, text width=2.7cm}]
    % (a) agente unico
    \node[lab] at (0,0) {(a) Agente\\único};
    \node[io] (ca) at (3.6,0) {Caso};
    \node[ag] (a1) at (6.2,0) {Avaliador};
    \node[io] (da) at (13.6,0) {Decisão};
    \draw[->] (ca) -- (a1); \draw[->] (a1) -- (da);
    % (b) pipeline
    \node[lab] at (0,-1.3) {(b) \textit{Pipeline}};
    \node[io] (cb) at (3.6,-1.3) {Caso};
    \node[ag] (ex) at (5.8,-1.3) {Extrator};
    \node[ag] (av) at (8.4,-1.3) {Avaliador};
    \node[ag] (si) at (11.4,-1.3) {Sintetizador};
    \node[io] (db) at (13.6,-1.3) {Decisão};
    \draw[->] (cb) -- (ex); \draw[->] (ex) -- node[above, font=\scriptsize]{perfil} (av);
    \draw[->] (av) -- node[above, font=\scriptsize]{parecer} (si); \draw[->] (si) -- (db);
    % (c) debate
    \node[lab] at (0,-3.0) {(c) Debate};
    \node[io] (cc) at (3.6,-3.0) {Caso};
    \node[ag] (pro) at (6.4,-2.4) {Debatedor pró};
    \node[ag] (con) at (6.4,-3.6) {Debatedor contra};
    \node[ag] (ju) at (10.4,-3.0) {Juiz};
    \node[io] (dc) at (13.6,-3.0) {Decisão};
    \draw[->] (cc) -- (pro.west); \draw[->] (cc) -- (con.west);
    \draw[<->] (pro) -- node[left, font=\scriptsize]{2 rodadas} (con);
    \draw[->] (pro) -- (ju); \draw[->] (con) -- (ju); \draw[->] (ju) -- (dc);
    % (d) hierarquica
    \node[lab] at (0,-5.0) {(d) Revisão hierárquica};
    \node[io] (cd) at (3.6,-5.0) {Caso};
    \node[ag] (ev) at (6.2,-5.0) {Avaliador};
    \node[ag] (rv) at (10.2,-5.0) {Revisor\\(manter/alterar)};
    \node[io] (dd) at (13.6,-5.0) {Decisão};
    \draw[->] (cd) -- (ev); \draw[->] (ev) -- node[above, font=\scriptsize]{parecer} (rv); \draw[->] (rv) -- (dd);
    \draw[->] (cd.south) to[bend right=12] node[below, font=\scriptsize]{caso original} (rv.south);
    % (e) ensemble
    \node[lab] at (0,-7.1) {(e) \textit{Ensemble}};
    \node[io] (ce) at (3.6,-7.1) {Caso};
    \node[ag, double copy shadow] (en) at (6.2,-7.1) {$K{=}5$ avaliadores\\independentes};
    \node[cd] (ag) at (10.2,-7.1) {Maioria (decisão)\\Média (pontuação)};
    \node[io] (de) at (13.6,-7.1) {Decisão};
    \draw[->] (ce) -- (en); \draw[->] (en) -- (ag); \draw[->] (ag) -- (de);
\end{tikzpicture}
\caption{As cinco topologias sob teste, todas com o mesmo LLM em todos os agentes (azul). Laranja: agregação determinística em código. No debate, a ordem de fala é sorteada e o juiz recebe o caso e a transcrição completa. Na revisão hierárquica, o revisor recebe o caso original e o parecer do avaliador. No \textit{pipeline}, cada agente recebe apenas a saída do agente anterior (o sintetizador também recebe o perfil extraído). O avaliador de (a), (d) e (e) usa o mesmo \textit{prompt}; (e) equivale a \textit{self-consistency} \citep{wang2023selfconsistency} aplicada a (a).}
\label{fig:topologias}
\end{figure}

\begin{table}[p]
\centering
\scriptsize
\caption{Especificação dos agentes. Entradas e saídas seguem esquemas JSON validados pelo orquestrador (Seção~\ref{sec:orquestracao}). ``Núcleo'' = instruções da tarefa + rubrica + esquema, idênticos para todos os agentes; descrição da vaga incluída no domínio de contratação.}
\label{tab:agentes}
\begin{tabularx}{\textwidth}{@{}p{1.4cm} p{1.6cm} L L L p{2.2cm}@{}}
\toprule
Topologia & Agente & Papel e escopo de decisão & Entrada & Saída & Término \\
\midrule
AU & Avaliador & Decide o caso inteiro sozinho. & Núcleo + caso & \texttt{\{justificativa, pontuacao, decisao\}} & 1 resposta válida \\
\midrule
PS & Extrator & Estrutura o caso: requisitos atendidos e experiência relevante, ou fatos, provas e teses. Não decide nem pontua. & Núcleo + caso & Perfil JSON com campos fixos e campo livre \texttt{outras\_informacoes} & 1 perfil válido \\
 & Avaliador & Avalia o perfil contra a rubrica; emite parecer preliminar. & Núcleo + perfil (sem o caso original) & \texttt{\{parecer, pontuacao\_prelim\}} & 1 parecer válido \\
 & Sintetizador & Emite a decisão final a partir do perfil e do parecer. & Núcleo + perfil + parecer & Esquema final & 1 resposta válida \\
\midrule
DB & Debatedor pró & Argumenta pelo desfecho favorável (avançar / absolver), sem decidir. & Núcleo + caso + transcrição até o turno & Argumento ($\le$ 250 \textit{tokens}) & 2 turnos \\
 & Debatedor contra & Argumenta pelo desfecho desfavorável (não avançar / condenar), sem decidir. & Idem & Idem & 2 turnos \\
 & Juiz & Decide com base no caso e nos argumentos; único agente com poder de decisão. & Núcleo + caso + transcrição completa & Esquema final & 1 resposta válida \\
\midrule
RH & Avaliador & Idêntico ao AU. & Núcleo + caso & Esquema final & 1 resposta válida \\ \addlinespace
 & Revisor & Supervisiona: verifica aderência do parecer à rubrica e pode manter ou alterar decisão e pontuação (veto/ajuste). & Núcleo + caso + saída do avaliador & \texttt{\{acao, justificativa, pontuacao, decisao\}} & 1 revisão (sem laço) \\
\midrule
EN & Avaliador$_k$, $k{=}1..5$ & Idêntico ao AU, executado 5 vezes de forma independente, sem visibilidade mútua. & Núcleo + caso & Esquema final & 1 resposta válida cada \\
\bottomrule
\end{tabularx}
\end{table}

\subsection{Hipóteses mecanicistas e predições mensuráveis}

A Tabela~\ref{tab:hipoteses} explicita, para cada topologia, o mecanismo pelo qual o viés poderia ser amplificado ou atenuado e a predição que o experimento pode refutar. As variáveis intermediárias citadas, como $\rho$, $\omega$, $\mu$ e $\pi_s$, são definidas na Seção~\ref{sec:metricas}.

\begin{table}[ht]
\centering
\footnotesize
\caption{Hipóteses por topologia e predições mensuráveis. $\Delta^{\mathrm{TF}}_\tau$ e $\Delta^{\CFR}_\tau$ são as diferenças entre a topologia $\tau$ e o agente único (Eq.~\ref{eq:contraste}). Há amplificação quando $|\Delta\mathrm{TF}_\tau|>|\Delta\mathrm{TF}_{\mathrm{AU}}|$ com o mesmo sinal.}
\label{tab:hipoteses}
\begin{tabularx}{\textwidth}{@{}p{1.6cm} L L@{}}
\toprule
Topologia & Mecanismo hipotetizado & Predição mensurável \\
\midrule
PS & \textbf{Filtragem ou saliência na extração.} O extrator pode descartar o atributo socioeconômico, irrelevante para a rubrica, o que atenuaria o viés. Ou pode registrá-lo em \texttt{outras\_informacoes}, tornando-o saliente para os estágios seguintes. & O sinal de $\Delta^{\mathrm{TF}}_{\mathrm{PS}}$ depende da taxa de retenção $\rho$, com atenuação se $\rho$ for baixa. Entre os pares em que o atributo é retido, $|\Delta\mathrm{TF}|$ e $\Delta\CFR$ são maiores do que entre os pares em que é descartado. \\
DB & \textbf{Conformidade e uso retórico.} Debatedores podem usar o atributo como argumento (``origem humilde'', ``defensor público''), e a dinâmica de conformidade pode levar a um consenso enviesado \citep{okawa2026emergence}. & Amplificação: $\Delta^{\mathrm{TF}}_{\mathrm{DB}}$ tem o mesmo sinal de $\Delta\mathrm{TF}_{\mathrm{AU}}$, e $|\SD_{\mathrm{DB}}|>|\SD_{\mathrm{AU}}|$. A taxa de menção ao atributo $\mu$ é maior nas transcrições do que nas justificativas do AU. \\
RH & \textbf{Ancoragem do revisor.} O revisor tende a confirmar o parecer recebido e herda o viés do avaliador em vez de corrigi-lo. & A taxa de alteração $\omega$ é baixa e não difere entre os níveis B e A. $\Delta^{\mathrm{TF}}_{\mathrm{RH}}$ é equivalente a zero (TOST, $\epsilon=0{,}05$). \\
EN & \textbf{Agregação reduz variância, não viés sistemático.} O voto majoritário de $K$ amostras mapeia a probabilidade individual $\pi$ em $P(\mathrm{Bin}(K,\pi)\ge 3)$, uma curva mais íngreme perto de $0{,}5$. Um deslocamento $\pi_{s,\mathrm{B}}-\pi_{s,\mathrm{A}}$ é, portanto, \emph{ampliado} em casos limítrofes. & $\CFR^{\mathrm{nulo}}_{\mathrm{EN}}<\CFR^{\mathrm{nulo}}_{\mathrm{AU}}$, isto é, menos ruído. $|\Delta\mathrm{TF}_{\mathrm{EN}}|\ge|\Delta\mathrm{TF}_{\mathrm{AU}}|$, com efeito concentrado nos casos com $\bar\pi_s\in[0{,}3;\,0{,}7]$ no AU. \\
\bottomrule
\end{tabularx}
\end{table}

\subsection{Raciocínio, planejamento e memória}

\paragraph{Raciocínio.} Todos os agentes decisores produzem uma justificativa curta \emph{antes} da pontuação e da decisão, uma vez que o esquema impõe essa ordem de campos. É uma cadeia de raciocínio leve e uniforme entre as topologias, o que evita que diferenças de profundidade de raciocínio se confundam com a topologia. \citet{morla2026agentfairbench} tratam o raciocínio elicitado como uma condição experimental própria, com controle de comprimento, o que indica que ele não deve variar junto com a topologia. Os debatedores raciocinam de forma argumentativa, conforme o papel atribuído, como no debate multiagente de \citet{du2024debate}. Aqui, porém, as posições são fixas (pró e contra) e um juiz separado decide.

\paragraph{Planejamento.} Nenhum agente planeja ou escolhe o próximo passo. O grafo de cada topologia é fixo e executado pelo orquestrador. A topologia é a variável independente do estudo e, por isso, não pode ser decidida dinamicamente pelos próprios agentes.

\paragraph{Memória.} Não há memória entre casos: cada caso, condição e réplica começa em contexto vazio, o que elimina contaminação e efeitos de ordem entre casos. Dentro de um caso, a memória é a \emph{visibilidade de mensagens} definida pela topologia:
\begin{itemize}
    \item \textbf{\textit{Pipeline}:} memória estritamente local, em que cada agente vê só a saída do anterior.
    \item \textbf{Debate:} transcrição compartilhada e acumulada.
    \item \textbf{Revisão hierárquica:} o revisor vê o caso e o parecer do avaliador.
    \item \textbf{\textit{Ensemble}:} isolamento total entre avaliadores.
\end{itemize}
Essa visibilidade é tratada como parte da topologia e documentada no Apêndice~A.1.

%%%%%%%%%%%%%%%%%%%%%%%%%%%%%%%%%%%%%%%%
\section{Orquestração e Uso de Ferramentas}
\label{sec:orquestracao}
%%%%%%%%%%%%%%%%%%%%%%%%%%%%%%%%%%%%%%%%

\paragraph{Padrão de orquestração.} Um orquestrador central determinístico, escrito em Python, implementa cada topologia como um grafo fixo de chamadas sobre a API compatível com OpenAI do vLLM. Frameworks como AutoGen \citep{wu2023autogen} ou CrewAI não são usados, pois injetam \textit{prompts} de sistema e lógicas de turno próprias, diferentes por padrão de colaboração. Isso introduziria diferenças de texto não controladas entre as topologias. Com o orquestrador próprio, a única diferença entre as topologias é a definida na Tabela~\ref{tab:agentes}.

\paragraph{Protocolo de comunicação.} Os agentes nunca se comunicam diretamente. Toda mensagem passa pelo orquestrador, que segue quatro regras:
\begin{itemize}
    \item valida a saída do agente contra o esquema JSON do papel;
    \item extrai apenas os campos previstos;
    \item monta a entrada do próximo agente a partir de um \textit{template} fixo;
    \item registra cada chamada num \textit{trace} com os campos \texttt{run\_id}, \texttt{case\_id}, condição, réplica, papel, \textit{prompt} completo, resposta bruta, resposta validada, \textit{tokens} de entrada e saída, latência e semente.
\end{itemize}
No debate, o orquestrador também faz o seguinte:
\begin{itemize}
    \item sorteia, com semente por caso e réplica, qual debatedor fala primeiro, controlando efeitos de ordem;
    \item alterna os turnos, de modo que cada debatedor vê os argumentos anteriores;
    \item encerra após 2 rodadas, totalizando 4 turnos, e passa a transcrição ao juiz.
\end{itemize}

\paragraph{Sementes e independência.} Cada chamada recebe uma semente derivada, por \textit{hash}, da tupla (caso, atributo, nível, réplica, papel).
\begin{itemize}
    \item \textbf{Sem números aleatórios comuns entre os níveis B e A.} Usar a mesma semente nos dois níveis reduziria artificialmente as inversões no par contrafactual em relação às réplicas, o que tornaria a CFR e o piso de ruído incomparáveis (Seção~\ref{sec:metricas}).
    \item \textbf{Não-determinismo residual.} O não-determinismo de \textit{batching} do vLLM não compromete o desenho, porque ele não depende de reprodução bit a bit das chamadas; as réplicas medem essa variabilidade.
\end{itemize}

\paragraph{Limites de execução.} Cada execução está sujeita aos seguintes limites:
\begin{itemize}
    \item \textbf{\textit{Tokens} de saída por chamada:} 400; 250 por turno de debate.
    \item \textbf{Rodadas:} 2 rodadas de debate e 1 revisão, sem laço de re-revisão.
    \item \textbf{\textit{Timeout}:} 120\,s por chamada.
    \item \textbf{Chamadas por decisão:} número fixo por topologia, igual a 1 (AU), 3 (PS), 5 (DB), 2 (RH) e 5 (EN).
\end{itemize}
Esses limites tornam o custo de cada topologia previsível e são registrados como covariáveis (Seção~\ref{sec:custo}).

\paragraph{Tratamento de falhas.} Uma saída que não passa no esquema é tratada assim:
\begin{enumerate}
    \item \textbf{Novas tentativas.} O orquestrador tenta novamente até 2 vezes. Cada nova tentativa acrescenta uma mensagem fixa (``Responda apenas no formato JSON especificado'') e usa uma nova semente derivada.
    \item \textbf{Decodificação restrita.} Persistindo a falha, faz uma última tentativa com decodificação restrita ao esquema (\textit{guided decoding} do vLLM).
    \item \textbf{Registro do estágio.} O estágio em que cada saída foi obtida (direto, após nova tentativa ou restrito) é registrado.
    \item \textbf{Saídas inválidas e recusas.} Saídas que continuam inválidas e recusas explícitas recebem o rótulo \texttt{invalido}. \emph{Nenhuma é descartada silenciosamente.}
\end{enumerate}
Falhas em agentes intermediários e no \textit{ensemble} seguem duas regras adicionais:
\begin{itemize}
    \item \textbf{Agentes intermediários.} Se o extrator, o avaliador do PS, um debatedor ou o avaliador do RH continuar inválido ou recusar a tarefa, a decisão inteira da topologia é marcada como \texttt{invalido}, pois a topologia não foi executada como especificada.
    \item \textbf{\textit{Ensemble}.} Um avaliador inválido é reamostrado com nova semente, até 2 vezes. Se continuar inválido, a decisão é marcada como \texttt{invalido}. Assim, toda decisão válida do \textit{ensemble} tem exatamente $K=5$ votos, e não há empate.
\end{itemize}
A taxa de \texttt{invalido} e a de recusa são comparadas entre os níveis B e A como métrica de viés própria. Uma recusa mais frequente para um dos grupos também é tratamento desigual. A análise principal exclui pares com qualquer lado inválido, e uma análise de sensibilidade imputa o desfecho desfavorável.

\paragraph{Tratamento de divergências.} As divergências entre agentes são resolvidas por regras fixas, nunca por um agente adicional:
\begin{itemize}
    \item \textbf{\textit{Ensemble}.} A decisão é o voto majoritário ($K=5$ é ímpar, logo não há empate), e a pontuação é a média.
    \item \textbf{Revisão hierárquica.} Prevalece a saída do revisor, que é sempre a final, mesmo com \texttt{acao=manter}. A incoerência entre \texttt{acao=manter} e uma decisão diferente da do avaliador é contada como erro de protocolo.
    \item \textbf{Debate.} Decide o juiz.
    \item \textbf{\textit{Pipeline}.} Decide o sintetizador.
    \item \textbf{Todas as topologias.} A incoerência entre pontuação e decisão, como pontuação acima de 50 com decisão desfavorável, é registrada como diagnóstico, mas não corrigida.
\end{itemize}

\paragraph{Ferramentas.} Os agentes sob teste \emph{não} têm ferramentas: não acessam busca, recuperação de documentos nem execução de código. Três razões justificam a escolha:
\begin{itemize}
    \item o caso completo cabe no \textit{prompt};
    \item resultados de ferramentas variam entre chamadas e topologias, adicionando uma fonte de variação confundida com a orquestração;
    \item a dimensão de uso de ferramentas já é estudada por \citet{morla2026agentfairbench}.
\end{itemize}
As ferramentas do sistema pertencem ao \textit{harness} e estão na Tabela~\ref{tab:ferramentas}.

\begin{table}[ht]
\centering
\footnotesize
\caption{Ferramentas do \textit{harness} (código determinístico; nenhuma é exposta aos agentes sob teste).}
\label{tab:ferramentas}
\begin{tabularx}{\textwidth}{@{}p{3.2cm} L L@{}}
\toprule
Ferramenta & Condição de uso & Validação da saída \\
\midrule
Renderizador / perturbador & Geração única dos pares, antes da execução & \textit{Diff} restrito ao \textit{slot}; \textit{hash} SHA-256 conferido a cada leitura \\
Validador de esquema JSON & Toda resposta de agente & Esquema por papel; campos obrigatórios; \texttt{decisao} em conjunto fechado; \texttt{pontuacao} inteira em $[0,100]$ \\
Detector de retenção/menção & Pós-processamento de \textit{traces} (PS, DB, RH) & Casamento exato e normalizado da forma de superfície do \textit{slot}; amostra de 100 casos conferida manualmente \\
Contador de \textit{tokens} e relógio & Toda chamada & \textit{Tokens} reportados pelo servidor conferidos com o \textit{tokenizer} local \\
Sonda de percepção & Validação dos léxicos (Seção~\ref{sec:dados}) & Concordância com a direção pretendida $\ge 80\%$ \\
\bottomrule
\end{tabularx}
\end{table}

\paragraph{Mecanismos de controle e intervenção humana.} Não há humano no laço de decisão, pois isso contaminaria a medida. Os pontos de intervenção ficam antes e depois da execução:
\begin{itemize}
    \item \textbf{H1, revisão dos casos-base e dos léxicos.} Ocorre antes do congelamento (Seção~\ref{sec:dados}).
    \item \textbf{H2, conferência da fidelidade dos pares.} O \textit{diff} automático cobre 100\% dos pares, e há leitura manual de 5\% deles.
    \item \textbf{H3, auditoria de \textit{traces}.} Uma amostra estratificada de 50 \textit{traces} por topologia e domínio é lida para verificar a aderência ao papel. O debatedor contra não pode, por exemplo, defender o candidato.
    \item \textbf{Pré-registro.} Hipóteses, métricas, testes e critérios de exclusão são congelados num documento versionado no repositório antes da execução principal, após o piloto.
    \item \textbf{Regra de parada.} A execução é interrompida se, nas primeiras 500 decisões de qualquer topologia, a taxa de \texttt{invalido} passar de 5\% ou a de erros de protocolo em H3 passar de 10\%. Nesse caso, os \textit{prompts} de papel são corrigidos, a alteração é documentada e a execução daquela topologia é reiniciada do zero.
\end{itemize}

%%%%%%%%%%%%%%%%%%%%%%%%%%%%%%%%%%%%%%%%
\section{Procedimentos de Avaliação}
\label{sec:avaliacao}
%%%%%%%%%%%%%%%%%%%%%%%%%%%%%%%%%%%%%%%%

\subsection{Métricas}
\label{sec:metricas}

\paragraph{Notação.} A notação usada nas métricas é a seguinte:
\begin{itemize}
    \item $d\in\{H,C\}$ é o domínio;
    \item $p$ é o atributo perturbado;
    \item $s\in S_{d}$ é o caso-base, com $|S_d|=N=200$;
    \item $\ell\in\{\mathrm{B},\mathrm{A}\}$ é o nível do atributo;
    \item $r\in\{1,\dots,R\}$, com $R=5$, é a réplica;
    \item $\tau\in\{\mathrm{AU},\mathrm{PS},\mathrm{DB},\mathrm{RH},\mathrm{EN}\}$ é a topologia;
    \item $Y^{\tau}_{s,\ell,r}\in\{0,1\}$ é a decisão final, com $1$ indicando o desfecho favorável;
    \item $X^{\tau}_{s,\ell,r}\in[0,100]$ é a pontuação final, orientada como na Seção~\ref{sec:dados}.
\end{itemize}
Para simplificar, omitem-se $d$ e $p$: toda métrica é calculada separadamente por (topologia, domínio, atributo).

\paragraph{CFR contrafactual e piso de ruído.} Adota-se a forma pareada da CFR de \citet{morla2026agentfairbench}, agregada sobre as réplicas:
\begin{equation}
\CFR^{\mathrm{pert}}_\tau=\frac{1}{NR}\sum_{s}\sum_{r}\ind\!\left[Y^\tau_{s,\mathrm{B},r}\neq Y^\tau_{s,\mathrm{A},r}\right].
\end{equation}
O piso de ruído é a mesma estatística calculada entre duas réplicas de uma entrada \emph{idêntica}, com pareamento cíclico $r\oplus1=(r \bmod R)+1$:
\begin{equation}
\CFR^{\mathrm{nulo}}_\tau=\frac{1}{2NR}\sum_{s}\sum_{\ell}\sum_{r}\ind\!\left[Y^\tau_{s,\ell,r}\neq Y^\tau_{s,\ell,r\oplus1}\right].
\end{equation}
As duas estatísticas comparam exatamente duas decisões, ou seja, têm a mesma aridade. Isso evita a inflação por aridade que \citet{morla2026agentfairbench} identificam ao comparar dispersões de $k>2$ grupos com um piso de duas execuções.

\paragraph{CFR em excesso.} É definida como
\begin{equation}
\Delta\CFR_\tau=\CFR^{\mathrm{pert}}_\tau-\CFR^{\mathrm{nulo}}_\tau.
\label{eq:dcfr}
\end{equation}

\paragraph{Diferença assinada de taxa de desfecho favorável.} É definida como
\begin{equation}
\Delta\mathrm{TF}_\tau=\frac{1}{NR}\sum_s\sum_r\left(Y^\tau_{s,\mathrm{B},r}-Y^\tau_{s,\mathrm{A},r}\right).
\label{eq:dtf}
\end{equation}
Valores negativos indicam desvantagem do nível socioeconômico baixo.

\paragraph{Relação entre as duas.} Sejam $\pi_{s,\ell}=P(Y^\tau_{s,\ell,r}=1)$ e $\Delta\pi_s=\pi_{s,\mathrm{B}}-\pi_{s,\mathrm{A}}$, e suponha réplicas independentes. Então $\mathbb{E}[\ind(Y_{s,\mathrm{B}}\neq Y_{s,\mathrm{A}})]=\pi_{s,\mathrm{B}}(1-\pi_{s,\mathrm{A}})+\pi_{s,\mathrm{A}}(1-\pi_{s,\mathrm{B}})$, e o piso médio vale $\pi_{s,\mathrm{B}}(1-\pi_{s,\mathrm{B}})+\pi_{s,\mathrm{A}}(1-\pi_{s,\mathrm{A}})$. Subtraindo,
\begin{equation}
\mathbb{E}\left[\Delta\CFR_\tau\right]=\frac{1}{N}\sum_s\Delta\pi_s^2=\overline{\Delta\pi}^{\,2}+\mathrm{Var}_s(\Delta\pi_s),
\qquad
\mathbb{E}\left[\Delta\mathrm{TF}_\tau\right]=\overline{\Delta\pi}.
\end{equation}
Isso tem quatro consequências:
\begin{itemize}
    \item \textbf{Robustez ao ruído.} As duas estatísticas medem o deslocamento da probabilidade de decisão causado \emph{apenas} pelo atributo. Nenhuma delas é afetada pelo fato de algumas topologias serem mais ou menos ruidosas. Isso é essencial, porque o voto majoritário do \textit{ensemble} reduz mecanicamente a CFR bruta, mesmo sem nenhuma alteração no viés.
    \item \textbf{$\Delta\CFR$ não tem sinal.} Ela é nula em expectativa se e somente se o atributo não afeta nenhum caso. Também capta efeitos heterogêneos, em que o atributo favorece o nível B em alguns casos e o A em outros, e que se cancelam em $\Delta\mathrm{TF}$.
    \item \textbf{$\Delta\CFR$ é de segunda ordem.} Um deslocamento sistemático de 5 pontos percentuais contribui com apenas $0{,}0025$ para $\Delta\CFR$, abaixo do que o desenho consegue distinguir (Seção~\ref{sec:analise}).
    \item \textbf{$\Delta\mathrm{TF}$ é de primeira ordem e tem sinal.} Recupera a direção do viés, o que é necessário porque a literatura registra inversões de direção entre gerações de modelos \citep{gao2026hiring}.
\end{itemize}

\paragraph{Hierarquia das métricas.} O desfecho principal continua sendo a decisão, como na Fase 1, e as duas estatísticas têm papéis distintos:
\begin{itemize}
    \item $\Delta\mathrm{TF}$ é a estatística do \textbf{teste confirmatório}.
    \item $\Delta\CFR$ é a \textbf{métrica principal de reporte}, presente em todas as tabelas, e o teste secundário de efeitos heterogêneos.
\end{itemize}

\paragraph{Contrastes de topologia.} O efeito da topologia é a diferença em relação à linha de base:
\begin{equation}
\Delta^{\mathrm{TF}}_\tau=\Delta\mathrm{TF}_\tau-\Delta\mathrm{TF}_{\mathrm{AU}},\qquad
\Delta^{\CFR}_\tau=\Delta\CFR_\tau-\Delta\CFR_{\mathrm{AU}},\qquad \tau\in\{\mathrm{PS},\mathrm{DB},\mathrm{RH},\mathrm{EN}\}.
\label{eq:contraste}
\end{equation}
A leitura de $\Delta^{\mathrm{TF}}_\tau$ depende de $\Delta\mathrm{TF}_{\mathrm{AU}}$:
\begin{itemize}
    \item há \emph{amplificação} se $\Delta^{\mathrm{TF}}_\tau$ tiver o mesmo sinal de $\Delta\mathrm{TF}_{\mathrm{AU}}$;
    \item há \emph{atenuação} se tiver sinal oposto sem ultrapassá-lo;
    \item há \emph{viés emergente} se $\Delta\mathrm{TF}_{\mathrm{AU}}$ for nulo e $\Delta\mathrm{TF}_\tau$ não for.
\end{itemize}

\paragraph{Métricas secundárias de pontuação.} São duas:
\begin{itemize}
    \item \textbf{MASD.} Para pares ($k=2$), a amplitude da definição de \citet{morla2026agentfairbench} é a diferença absoluta:
    \begin{equation}
    \MASD_\tau=\frac{1}{NR}\sum_s\sum_r\left|X^\tau_{s,\mathrm{B},r}-X^\tau_{s,\mathrm{A},r}\right| .
    \end{equation}
    Ela é reportada junto ao seu piso, calculado sobre réplicas como na CFR, e ao excesso $\Delta\MASD_\tau$.
    \item \textbf{Diferença assinada de pontuação.} É o análogo de $\Delta\mathrm{TF}$ na pontuação, que capta deslocamentos que não chegam a cruzar o limiar de decisão:
    \begin{equation}
    \SD_\tau=\frac{1}{NR}\sum_s\sum_r\left(X^\tau_{s,\mathrm{B},r}-X^\tau_{s,\mathrm{A},r}\right).
    \end{equation}
\end{itemize}

\paragraph{Variáveis intermediárias (QP3).} Cada variável está associada a uma das hipóteses da Tabela~\ref{tab:hipoteses}:
\begin{itemize}
    \item $\rho$ é a taxa de retenção, a fração de saídas do extrator (PS) que contém a forma de superfície do atributo.
    \item $\omega_\ell$ é a taxa de alteração, a fração de casos em que o revisor (RH) muda a decisão do avaliador, por nível $\ell$.
    \item $\mu$ é a taxa de menção ao atributo nas transcrições do debate (DB) e nas justificativas das demais topologias.
    \item $\bar\pi_s$ é a probabilidade média de desfecho favorável do caso $s$ no AU. Ela é usada para estratificar casos limítrofes.
    \item Concordância entre avaliadores do \textit{ensemble}: kappa de Fleiss.
\end{itemize}

\paragraph{Métricas de validade e operacionais.} Além das métricas de viés, são registradas:
\begin{itemize}
    \item a taxa de \texttt{invalido} e a de recusa por nível, comparadas entre B e A pelo teste de McNemar sobre os pares;
    \item a taxa de erros de protocolo (H3);
    \item a taxa de incoerência entre pontuação e decisão;
    \item chamadas por decisão, \textit{tokens} de entrada e de saída por decisão e latência por decisão (mediana e p95).
\end{itemize}
Não há métrica de acurácia, pois os casos não têm rótulo verdadeiro. A sensibilidade a fatores legítimos é verificada pelo controle positivo descrito abaixo. Também não há avaliação com usuários-alvo, como recrutadores ou jurados, porque a Fase 1 excluiu testes com usuários reais do escopo. O benchmark se destina a equipes de engenharia e a auditores, que usam diretamente as métricas acima.

\subsection{Protocolo experimental}
\label{sec:protocolo}

\paragraph{Desenho.} O experimento é fatorial completo e intra-casos, isto é, todo caso passa por todas as condições:
\[
2 \text{ domínios} \times 3 \text{ atributos} \times 2 \text{ níveis} \times 5 \text{ topologias} \times R{=}5 \text{ réplicas} \times N{=}200 \text{ casos}.
\]
Isso resulta em $60.000$ decisões finais por modelo, executadas primeiro com M1 e depois repetidas integralmente com M2. Como cada caso passa por todas as topologias, os contrastes da Eq.~\ref{eq:contraste} são pareados por caso, o que remove a variância entre casos.

\paragraph{\textit{Baselines} e controles.}
\begin{itemize}
    \item \textbf{Linha de base de topologia:} o agente único (AU), com o mesmo \textit{prompt} do avaliador de RH e EN.
    \item \textbf{Linha de base de ruído:} as réplicas idênticas, ou seja, $\CFR^{\mathrm{nulo}}$ para cada topologia.
    \item \textbf{Controle positivo de sensibilidade.} Em 50 casos por domínio, o nível de qualificação ou a força da prova é deslocado em um estrato, mantendo o resto constante. Espera-se $|\Delta\mathrm{TF}|$ e $\Delta\CFR$ muito maiores do que os de qualquer atributo socioeconômico. Se isso não ocorrer, a topologia não está lendo o caso, e os resultados de viés não são interpretáveis.
    \item \textbf{Controle positivo de detecção.} Em 50 casos por domínio, o AU recebe uma instrução explícita para favorecer o nível A. O \textit{pipeline} de análise precisa detectar $\Delta\mathrm{TF}<0$ e $\SD<0$ com $p<0{,}05$, o que demonstra que o instrumento recupera um viés conhecido \citep{morla2026agentfairbench}.
    \item \textbf{Sensibilidade à temperatura.} O AU é executado com $T=0$ em todos os pares, para verificar se a direção de $\Delta\mathrm{TF}$ e de $\SD$ se mantém.
\end{itemize}

\paragraph{Condições controladas e fatores de confusão.} Os principais fatores de confusão são tratados assim:
\begin{itemize}
    \item \textbf{Modelo, \textit{prompts} e decodificação.} São fixados pelos invariantes da Seção~\ref{sec:arquitetura}.
    \item \textbf{Ordem de fala no debate.} É sorteada.
    \item \textbf{Ordem de execução das chamadas.} É aleatorizada entre condições, para que a deriva do servidor não se correlacione com o nível.
    \item \textbf{Comprimento do texto.} A diferença entre os níveis B e A é inferior a 5 \textit{tokens}.
    \item \textbf{Número de chamadas e de \textit{tokens}.} Difere entre topologias por definição e \emph{não} é equalizado, pois faz parte do que distingue cada arquitetura. Por isso, é reportado como covariável (Seção~\ref{sec:custo}), e as conclusões são formuladas como ``a topologia, com seu custo característico, altera o viés''.
\end{itemize}

\subsection{Análise estatística}
\label{sec:analise}

\paragraph{Unidade de análise.} A unidade de análise é o caso-base $s$. Réplicas e níveis são medidas repetidas do mesmo caso, e tratá-los como independentes estreitaria artificialmente os intervalos.

\paragraph{Intervalos de confiança.} Para cada métrica, os intervalos de 95\% são obtidos por \textit{bootstrap} por conglomerado: os casos são reamostrados com reposição, $B=10.000$ vezes, com intervalo percentil. Para taxas simples, usa-se o intervalo de Wilson.

\paragraph{Testes confirmatórios.} Para cada topologia $\tau$ e domínio $d$, calcula-se a diferença por caso, agregada sobre os três atributos do domínio, que estão todos orientados de B para A:
\[
\delta_s=\frac{1}{3}\sum_{p}\left(\Delta\mathrm{TF}_{\tau,s,p}-\Delta\mathrm{TF}_{\mathrm{AU},s,p}\right).
\]
Sobre $\{\delta_s\}$ aplica-se um teste de permutação por troca de sinal, com $10.000$ permutações e hipótese bilateral.
\begin{itemize}
    \item A família confirmatória tem $4$ topologias $\times$ $2$ domínios, ou seja, $8$ testes.
    \item Os testes são corrigidos por Holm \citep{holm1979} para uma taxa de erro familiar de $0{,}05$.
    \item O teste é executado sobre o modelo M1. O modelo M2 é tratado como replicação, e um efeito é considerado replicado se tiver o mesmo sinal e $p<0{,}05$ não corrigido em M2.
\end{itemize}
A agregação sobre os atributos é o que torna a família pequena e o teste poderoso. O detalhamento por atributo é secundário.

\paragraph{Testes secundários.} São controlados por Benjamini--Hochberg com $q=0{,}05$:
\begin{itemize}
    \item os $24$ contrastes $\Delta^{\mathrm{TF}}_\tau$ por atributo;
    \item os contrastes $\Delta^{\CFR}_\tau$;
    \item os testes de $\Delta\CFR_\tau>0$ e de $\SD_\tau\neq0$ (Wilcoxon pareado) para cada topologia isolada;
    \item as comparações entre B e A nas taxas de inválido e de recusa (McNemar);
    \item as predições da Tabela~\ref{tab:hipoteses}.
\end{itemize}
Como análise de robustez, ajusta-se um modelo logístico de efeitos mistos para a decisão:
\[
\mathrm{logit}\,P(Y=1)\sim \tau\times d\times p\times \ell + (1\mid s).
\]
O efeito da topologia sobre o viés corresponde aos termos de interação $\tau\times\ell$.

\paragraph{Ausência de efeito.} A pergunta ``a topologia importa?'' admite resposta negativa, que precisa ser demonstrada, e não apenas presumida a partir de $p>0{,}05$. Para isso, testa-se a equivalência de cada um dos 8 contrastes confirmatórios por dois testes unilaterais (TOST, $\alpha=0{,}05$) com margem $\epsilon=0{,}05$, isto é, 5 pontos percentuais na probabilidade de desfecho favorável.

Essa margem tem uma limitação que precisa ser explícita. Disparidades da ordem de 1 a 3 pontos percentuais, como as reportadas por \citet{an2025resume} para LLM único em contratação, ficam abaixo da resolução deste desenho. Uma conclusão de equivalência significa ``nenhuma diferença maior que 5 pontos percentuais'', e não ``nenhuma diferença''.

\paragraph{Poder e tamanho amostral.} O poder foi estimado por simulação de Monte Carlo, com $1.000$ repetições, sob as seguintes hipóteses:
\begin{itemize}
    \item $\pi_s\sim\mathcal{U}[0{,}1;\,0{,}9]$, coerente com a seleção piloto;
    \item decisões de Bernoulli independentes entre réplicas;
    \item AU sem viés e deslocamento homogêneo $\Delta\pi$ na topologia testada;
    \item $N=200$ e $R=5$.
\end{itemize}
Os resultados foram:
\begin{itemize}
    \item \textbf{Teste confirmatório.} O desvio-padrão das diferenças por caso agregadas é $\sigma_\delta\approx0{,}23$. Com Holm em 8 testes, o efeito mínimo detectável com poder de 80\% é $\mathrm{MDE}\approx(z_{1-0{,}05/16}+z_{0{,}8})\,\sigma_\delta/\sqrt{N}\approx0{,}06$. O poder simulado é $0{,}63$ para $\Delta\pi=0{,}05$ e $\approx1{,}00$ para $\Delta\pi=0{,}10$.
    \item \textbf{Equivalência.} Sob contraste nulo verdadeiro, o TOST com $\epsilon=0{,}05$ declara equivalência em 85\% das repetições.
    \item \textbf{Contrastes por atributo.} Com 24 testes, o poder é $0{,}69$ para $\Delta\pi=0{,}10$.
    \item \textbf{$\Delta^{\CFR}_\tau$ como teste confirmatório.} O poder seria $\le0{,}04$ para $\Delta\pi=0{,}10$ e só chegaria a $0{,}73$ com $\Delta\pi=0{,}30$ (24 testes). Este é o motivo do ajuste de escopo 4.
\end{itemize}
A simulação ignora variância adicional entre agentes e casos. Por isso, $\sigma_\delta$ será reestimado no piloto, com 40 casos por domínio e todas as topologias. Se o MDE resultante passar de $0{,}08$, o desenho é ampliado antes do pré-registro, de uma de duas formas:
\begin{itemize}
    \item $R$ passa para $10$, o que reduz $\sigma_\delta$ por $\approx\sqrt2$ na simulação e dobra o custo da Tabela~\ref{tab:custo};
    \item ou $N$ passa para $400$.
\end{itemize}
% TODO: atualizar sigma_delta e MDE com os dados do piloto (script de simulação a ser versionado no repositório).

\subsection{Custo e latência}
\label{sec:custo}

A Tabela~\ref{tab:custo} estima o volume de chamadas e \textit{tokens}. A hipótese é de, em média, $1.500$ \textit{tokens} de entrada e $300$ de saída por chamada; as chamadas de juiz e de revisor são maiores, as de debatedor menores. O custo monetário é nulo por \textit{token}, porque a inferência é local, mas o custo em GPU-hora é registrado. O custo efetivo por decisão e a latência (mediana e p95) de cada topologia são resultados reportados ao lado do viés. Isso permite à equipe que adota uma arquitetura ponderar viés contra custo, na linha da análise de custo por orquestração de \citet{sayedali2026harness}.

\begin{table}[ht]
\centering
\footnotesize
\caption{Orçamento estimado por modelo (M1 ou M2). Decisões por topologia: $2\times3\times2\times5\times200=12.000$.}
\label{tab:custo}
\begin{tabular}{@{}lrrr@{}}
\toprule
Topologia & Chamadas/decisão & Chamadas & \textit{Tokens} (estim., $10^6$) \\
\midrule
AU & 1 & 12.000 & 22 \\
PS & 3 & 36.000 & 65 \\
DB & 5 & 60.000 & 108 \\
RH & 2 & 24.000 & 43 \\
EN & 5 & 60.000 & 108 \\
\midrule
Total experimento principal & 16 & 192.000 & $\approx$ 346 \\
Piloto, controles e sonda & -- & $\approx$ 64.000 & $\approx$ 115 \\
\bottomrule
\end{tabular}
\end{table}
% TODO: medir throughput real no piloto e converter em GPU-horas; confirmar acesso a GPU (cluster do IC ou alternativa).

\subsection{Critérios de sucesso}

O sucesso é avaliado em dois níveis.

\paragraph{Validade do instrumento.} Estes critérios são pré-requisitos e, se algum falhar, os resultados de viés não são reportados como conclusivos:
\begin{itemize}
    \item 100\% dos pares passam no \textit{diff};
    \item ao menos 80\% das variantes de léxico passam na sonda de percepção;
    \item os dois controles positivos são detectados;
    \item a taxa de \texttt{invalido} fica abaixo de 2\% por topologia após as novas tentativas;
    \item os erros de protocolo em H3 ficam abaixo de 10\%.
\end{itemize}

\paragraph{Resposta científica.} As questões de pesquisa são respondidas com um entregável bem definido:
\begin{itemize}
    \item \textbf{QP1:} cada um dos 8 contrastes confirmatórios recebe uma de três conclusões: diferente de zero (Holm), equivalente a zero com margem de 5 pontos percentuais (TOST) ou inconclusivo.
    \item \textbf{QP2:} consistência de sinal e de significância entre domínios, atributos e modelos.
    \item \textbf{QP3:} confirmação ou refutação de cada predição da Tabela~\ref{tab:hipoteses}.
\end{itemize}
O projeto é bem-sucedido se entregar essa resposta com o instrumento validado, independentemente da direção do resultado. Um resultado nulo bem estimado é uma resposta válida à pergunta do título.

\subsection{Ameaças à validade e limitações}
\label{sec:ameacas}

\paragraph{Validade de construto.} Há três ameaças principais:
\begin{itemize}
    \item \textbf{Atributos como \textit{proxies}.} Os atributos operacionalizam status socioeconômico por \textit{proxies}, que também podem sinalizar raça, região ou outros fatores. Bairros de periferia em São Paulo, por exemplo, têm composição racial distinta da dos bairros centrais. A sonda de percepção mitiga, mas não elimina, essa sobreposição, e os resultados são formulados como resposta ao \textit{proxy}, não à classe em abstrato.
    \item \textbf{Ocupação do réu.} Pode ser lida como informação de contexto e não apenas de status.
    \item \textbf{Decisão binária forçada.} Simplifica decisões reais, que incluem graus e alternativas, como a desclassificação do crime ou o pedido de mais informações.
\end{itemize}

\paragraph{Validade interna.} Há quatro ameaças principais:
\begin{itemize}
    \item \textbf{Número de chamadas e de \textit{tokens}.} Covaria com a topologia por construção. Uma diferença entre EN e AU, por exemplo, não distingue ``agregação'' de ``mais computação''. Isso é registrado, mas não resolvido.
    \item \textbf{Seleção dos casos pelo AU.} A seleção de casos limítrofes pelo comportamento do AU pode favorecer a detecção de efeitos na linha de base. Para verificar isso, os contrastes da Eq.~\ref{eq:contraste} são reportados também por estrato de $\bar\pi_s$.
    \item \textbf{Texto dos sufixos de papel.} É inevitavelmente diferente entre topologias, mas é curto e é publicado.
    \item \textbf{Não-determinismo do servidor.} É absorvido pelas réplicas, mas impede a reprodução exata das chamadas; a reprodução é garantida a partir dos \textit{traces}.
\end{itemize}

\paragraph{Validade externa.} Há quatro ameaças principais:
\begin{itemize}
    \item \textbf{Escala dos modelos.} Dois modelos abertos de 8B parâmetros não representam modelos comerciais de fronteira. A direção do viés de LLMs em contratação muda entre gerações \citep{gao2026hiring}.
    \item \textbf{Casos sintéticos.} São mais curtos e limpos que currículos e autos reais.
    \item \textbf{Contexto e idioma.} O contexto é paulistano e em português, o que limita a generalização para outros contextos e idiomas, embora seja justamente o contexto de interesse do projeto.
    \item \textbf{Implementação das topologias.} Cada topologia tem uma única implementação, com número de rodadas, $K$ e papéis fixos. Outras parametrizações podem se comportar de forma diferente, e o estudo não varia esses parâmetros.
\end{itemize}

\paragraph{Validade de conclusão.} Há duas ameaças principais:
\begin{itemize}
    \item \textbf{Graus de liberdade do pesquisador.} A família de 8 testes confirmatórios, a margem de equivalência e as decisões de exclusão são fixadas no pré-registro, o que reduz esses graus de liberdade.
    \item \textbf{Poder estatístico.} O poder simulado supõe ruído apenas binomial, e o $\sigma_\delta$ real pode ser maior; o piloto existe para estimá-lo antes da execução principal. Mesmo no cenário simulado, efeitos menores que 5 pontos percentuais não são detectáveis.
\end{itemize}

%%%%%%%%%%%%%%%%%%%%%%%%%%%%%%%%%%%%%%%%
\section*{Apêndice}
\label{sec:apendice}
%%%%%%%%%%%%%%%%%%%%%%%%%%%%%%%%%%%%%%%%

\subsection*{A.1 Núcleo de \textit{prompt} e esquemas}
\phantomsection

Núcleo comum a todos os agentes, ilustrado para o domínio de contratação; a versão final de todos os \textit{prompts} será publicada no repositório:

\begin{small}
\begin{verbatim}
Você está avaliando uma candidatura para a vaga descrita abaixo.
Use apenas os critérios da rubrica: requisitos obrigatórios,
requisitos desejáveis e experiência relevante.
[VAGA]  {descricao_da_vaga}
[CASO]  {texto_do_caso | perfil | transcrição, conforme o papel}
Responda somente em JSON:
{"justificativa": "<até 80 palavras>",
 "pontuacao": <inteiro de 0 a 100>,
 "decisao": "avancar" | "nao_avancar"}
\end{verbatim}
\end{small}

Sufixos de papel (resumo):
\begin{itemize}
    \item \textbf{Extrator:} ``estruture o caso nos campos dados; não avalie''.
    \item \textbf{Avaliador do PS:} ``emita parecer preliminar sobre o perfil''.
    \item \textbf{Sintetizador:} ``decida com base no perfil e no parecer''.
    \item \textbf{Debatedor pró/contra:} ``argumente que a decisão deve ser favorável/desfavorável; não decida''.
    \item \textbf{Juiz:} ``decida com base no caso e nos argumentos''.
    \item \textbf{Revisor:} ``verifique o parecer contra a rubrica e mantenha ou altere a decisão, justificando''.
\end{itemize}
No domínio penal, o núcleo substitui a vaga pelas instruções ao jurado (decidir sobre autoria e materialidade com base nas provas apresentadas) e usa \texttt{"decisao": "condenar" | "absolver"} e \texttt{"conviccao\_culpa"} em lugar de \texttt{"pontuacao"}.

\subsection*{A.2 Léxicos de perturbação}
\phantomsection

Cada nível tem cinco variantes de superfície. Os exemplos estão na Tabela~\ref{tab:perturbacoes}, e a lista completa será congelada após a revisão H1 e a sonda de percepção.
% TODO: completar as 5 variantes por nível após verificação no Mapa da Desigualdade 2024 e validação pela sonda.



%%%%%%%%%%%%%%%%%%%%%%%%%%%%%%%%%%
%          BIBLIOGRAFIA
%%%%%%%%%%%%%%%%%%%%%%%%%%%%%%%%%%
\newpage
\bibliographystyle{plainnat} % TODO: template original usa aasjournal.bst; pacote não instalado localmente (disponível via texlive-aastex). Confirmar com a disciplina se plainnat é aceito ou se aasjournal.bst é obrigatório antes da entrega.
\bibliography{../references}{}



\end{document}
